\renewcommand{\iint}{\mathop{\int\mkern-10mu\int}\nolimits}
\newcommand{\dint}[1]{\iint_{\mathcal{D}} #1 \, \dd{A}}

\section{Double Integrals}
\begin{definition}{Imprecise Definition of Double Integral}{def_double_integral_imprecise}
Let $f(x, y)$ be a function defined on a closed and bounded region $\mathcal{D}$ in the $xy$-plane.
The \textbf{double integral} is defined by
\[
    \dint{f(x, y)}
\]
where $\dd{A}$ is the area element in the $xy$-plane.

This integral can be interpreted as the volume under the surface $z = f(x, y)$ and above the region $\mathcal{D}$ in the $xy$-plane.
\end{definition}

Below are some properties of double integrals that will be useful in evaluating them.
\begin{itemize}
    \item $\displaystyle\dint{f(x, y) + g(x, y)} = \dint{f(x, y)} + \dint{g(x, y)}$ \\[-0.5em]
    \item $\displaystyle\dint{c \cdot f(x, y)} = c \cdot \dint{f(x, y)}$ \\[-0.5em]
    \item If $f(x, y) \leq g(x, y)$ for all $(x, y) \in \mathcal{D}$, then $\displaystyle\dint{f(x, y)} \leq \dint{g(x, y)}$.
\end{itemize}

\subsection{Riemann Sums and the Definition of Double Integrals}
Recall the concept of Riemann sums for single-variable functions.

The Riemann sum for a function $f(x)$ over an interval $[a, b]$ is given by
\[
    \sum_{i=1}^n f(p_i) \Delta x_i
\]
where $\Delta x_i$ is the width of the $i$-th subinterval,
and $p_i$ is a sample point in that subinterval, as indicated in the figure below.
\begin{center}
\begin{tikzpicture}[x=1.15cm,y=.9cm,font=\small,>=Stealth,
    declare function={g(\x)=1.25 + .65*sin(42 * \x);}]
    % Each rectangle meets the curve at its midpoint sample.
    \foreach \a/\b/\i in {.4/1.2/1,1.2/2.3/2,2.3/3.2/3,3.2/4.6/4}{
        \pgfmathsetmacro{\p}{(\a + \b)/2}
        \pgfmathsetmacro{\h}{g(\p)}
        \filldraw[fill=orange!45,fill opacity=.4,draw=orange!65!black]
            (\a,0) rectangle (\b,\h);
        \draw[dashed,black!55] (\p,0)--(\p,\h);
        \fill (\p,\h) circle (1.4pt);
        \node[below] at (\p,0) {$p_{\i}$};
    }
    \draw[->] (0,0)--(5.2,0) node[right] {$x$};
    \draw[->] (0,0)--(0,2.5) node[above] {$y$};
    \draw[blue!70!black,thick,domain=.25:4.8,samples=65]
        plot (\x,{g(\x)}) node[right] {$y = f(x)$};
    \node[below left] at (.4,0) {$a$};
    \node[below right] at (4.6,0) {$b$};
    \draw[<->] (1.2,-.65)--(2.3,-.65)
        node[midway,below] {$\Delta x_i$};
\end{tikzpicture}
\end{center}


We can extend this idea to functions of two variables.
Assume that $\mathcal{D} = [a, b] \times [c, d]$ is rectangular region in the $xy$-plane.

\begin{center}
\begin{tikzpicture}[font=\small,>=Stealth,
    declare function={
        f(\x,\y)=2.65 - .16 * (\x - 1.5)^2 - .12 * (\y - 1.5)^2;
    }]
    \tikzset{
        top face/.style={
            fill=orange!25,fill opacity=.35,
            draw=orange!45!black,line width=.45pt
        },
        side face/.style={
            fill=orange!45,fill opacity=.25,
            draw=orange!65!black,line width=.45pt
        },
        surface mesh/.style={
            draw=blue!85!black,draw opacity=.85,line width=.15pt
        }
    }

    % Both panels use exactly the same surface.
    \newcommand{\drawSurface}{
        \path[fill=cyan!40,fill opacity=.18]
            plot[domain=0:3,variable=\t,samples=25]
                (\t,0,{f(\t,0)})
            -- plot[domain=0:3,variable=\t,samples=25]
                (3,\t,{f(3,\t)})
            -- plot[domain=3:0,variable=\t,samples=25]
                (\t,3,{f(\t,3)})
            -- plot[domain=3:0,variable=\t,samples=25]
                (0,\t,{f(0,\t)}) -- cycle;
        \foreach \q in {0,.3,...,3}{
            \draw[surface mesh]
                plot[domain=0:3,variable=\t,samples=25]
                (\t,\q,{f(\t,\q)});
            \draw[surface mesh]
                plot[domain=0:3,variable=\t,samples=25]
                (\q,\t,{f(\q,\t)});
        }
        \node[above right,text=blue!75!black] at (0,3,3)
            {$z = f(x,y)$};
    }

    \foreach \panel in {0,1}{
        \begin{scope}[shift={({6.8 * \panel},0)},
            x={(-.64cm,-.42cm)},
            y={(.74cm,-.34cm)},
            z={(0cm,.8cm)}]

            % Base rectangle and coordinate axes.
            \filldraw[fill=orange!5,draw=black!45]
                (0,0,0)--(3,0,0)--(3,3,0)--(0,3,0)--cycle;
            \draw[->] (0,0,0)--(3.65,0,0)
                node[below left] {$x$};
            \draw[->] (0,0,0)--(0,3.65,0)
                node[right] {$y$};
            \draw[->] (0,0,0)--(0,0,3.65)
                node[above] {$z$};

            \ifnum\panel=0
                % A representative cell of a nonuniform partition.
                \foreach \q in {1.2,2.4}{
                    \draw[orange!40]
                        (\q,0,0)--(\q,3,0)
                        (0,\q,0)--(3,\q,0);
                }

                % The sample is the center of the rectangular base.
                \pgfmathsetmacro{\h}{f(1.8,1.8)}

                % Hidden edges.
                \draw[dashed,orange!50!black]
                    (1.2,1.2,\h)--(1.2,1.2,0)--(2.4,1.2,0)
                    (1.2,1.2,0)--(1.2,2.4,0);

                % Translucent side faces and horizontal top.
                \filldraw[side face]
                    (2.4,1.2,0)--(2.4,2.4,0)--
                    (2.4,2.4,\h)--(2.4,1.2,\h)--cycle;
                \filldraw[side face]
                    (1.2,2.4,0)--(2.4,2.4,0)--
                    (2.4,2.4,\h)--(1.2,2.4,\h)--cycle;
                \filldraw[top face]
                    (1.2,1.2,\h)--(2.4,1.2,\h)--
                    (2.4,2.4,\h)--(1.2,2.4,\h)--cycle;

                \draw[dashed,black!50]
                    (1.8,1.8,0)--(1.8,1.8,\h);

                % Overlay the surface so its intersection remains visible.
                \drawSurface
                \fill (1.8,1.8,\h) circle (1.6pt);
                \fill (1.8,1.8,0) circle (1.3pt);
                \draw (1.8,1.8,0)--(3.5,.1,.75)
                    node[left] {$P_{ij}$};

                % Dimension lines outside the base.
                \draw[densely dotted]
                    (1.2,2.4,0)--(1.2,3.6,0)
                    (2.4,2.4,0)--(2.4,3.6,0)
                    (2.4,1.2,0)--(3.6,1.2,0)
                    (2.4,2.4,0)--(3.6,2.4,0);
                \draw[<->] (1.2,3.6,0)--(2.4,3.6,0)
                    node[midway,below right] {$\Delta x_i$};
                \draw[<->] (3.6,1.2,0)--(3.6,2.4,0)
                    node[midway,below left] {$\Delta y_j$};

            \else
                % Draw exposed column faces from back to front.
                % Interior walls are omitted to keep the surface readable.
                \foreach \s in {0,...,8}{
                    \foreach \i in {0,...,4}{
                        \pgfmathtruncatemacro{\j}{\s - \i}
                        \ifnum\j<0\else\ifnum\j>4\else
                            \pgfmathsetmacro{\u}{.6 * \i}
                            \pgfmathsetmacro{\v}{.6 * \j}

                            % Each box height is evaluated at its midpoint.
                            \pgfmathsetmacro{\h}{f(\u + .3,\v + .3)}
                            \pgfmathsetmacro{\hx}{
                                \i < 4 ? min(\h,f(\u + .9,\v + .3)) : 0
                            }
                            \pgfmathsetmacro{\hy}{
                                \j < 4 ? min(\h,f(\u + .3,\v + .9)) : 0
                            }

                            \filldraw[side face]
                                ({\u + .6},\v,\hx)--
                                ({\u + .6},{\v + .6},\hx)--
                                ({\u + .6},{\v + .6},\h)--
                                ({\u + .6},\v,\h)--cycle;
                            \filldraw[side face]
                                (\u,{\v + .6},\hy)--
                                ({\u + .6},{\v + .6},\hy)--
                                ({\u + .6},{\v + .6},\h)--
                                (\u,{\v + .6},\h)--cycle;
                            \filldraw[top face]
                                (\u,\v,\h)--({\u + .6},\v,\h)--
                                ({\u + .6},{\v + .6},\h)--
                                (\u,{\v + .6},\h)--cycle;
                        \fi\fi
                    }
                }

                % Surface and exact sample intersections.
                \drawSurface
                \foreach \i in {0,...,4}{
                    \foreach \j in {0,...,4}{
                        \pgfmathsetmacro{\px}{.6 * \i + .3}
                        \pgfmathsetmacro{\py}{.6 * \j + .3}
                        \fill (\px,\py,{f(\px,\py)}) circle (.5pt);
                    }
                }
            \fi
        \end{scope}
    }
\end{tikzpicture}
\end{center}


By this, we could write a precise definition of the double integral as follows:
\begin{definition}{Precise Definition of Double Integral}{def_double_integral_precise}
Let $f(x, y)$ be a function defined on a closed and bounded region $\mathcal{D}$ in the $xy$-plane.

We partition $\mathcal{D}$ into $m$ subintervals in the $x$-direction and $n$ subintervals in the $y$-direction, forming a grid of rectangles.

Let $\Delta x_i$ and $\Delta y_j$ be the widths of the subintervals in the $x$ and $y$ directions, respectively,
and let $P_{ij}$ be a sample point in the rectangle formed by the $i$-th subinterval in the $x$-direction and the $j$-th subinterval in the $y$-direction.

The \textbf{double integral} of $f(x, y)$ over the region $\mathcal{D}$ is defined as the limit of the Riemann sum as the maximum width of the subintervals approaches zero:
\[
    \dint{f(x, y)} = \lim_{\norm{P} \to 0} \sum_{i=1}^m \sum_{j=1}^n f(P_{ij}) \Delta A_{ij}
\]
\end{definition}

Notice that the $\Delta A_{ij}$ here does not always exist. If it does, we call the function $f(x, y)$ \textbf{integrable} over the region $\mathcal{D}$.
Also, if a function is continuous over a closed and bounded region, then it is integrable over that region.

\subsection{Iterated Integrals and Fubini's Theorem}
To evaluate double integrals, we can use Fubini's Theorem, which allows us to compute the double integral as an iterated integral.
\begin{theorem}{Fubini's Theorem}{thm_fubini}
\[
    \dint{f(x, y)} = \int_a^b \int_c^d f(x, y) \, \dd{y} \, \dd{x} = \int_c^d \int_a^b f(x, y) \, \dd{x} \, \dd{y}
\]
\end{theorem}